

\section{The \textbf{\textit{n}}{\fontfamily{ptm}\selectfont DNA} Cartograph: Latent Semantic Genome of Foundation Models}
\label{sec:ndna_score}

\noindent
\textbf{Before we unveil \textit{\emph{n}DNA}, we must confront a foundational question:} \emph{What qualifies as heritability in artificial cognition}? Conventional artifacts--\textit{weights}, \textit{activations}, or the output \textit{behavior}--are mere \textbf{epiphenomena of training}. In contrast, \textbf{\textit{n}}\textnormal{DNA} seeks to capture a model’s \emph{semantic genome}: the \textit{latent organizational structures} that govern how knowledge is internally \emph{represented}, \emph{adapted}, and \emph{transmitted} across fine-tuning, distillation, pruning, and deployment. To chart the \textbf{semantic ancestry} of AI systems, we must move beyond output-level metrics and embrace a deeper epistemic foundation--one that traces not just what models \emph{say}, but how they \emph{reason}, \emph{evolve}, and \emph{remember}. We argue that \textbf{\textit{n}}\textnormal{DNA} constitutes this missing genomic trace: a \textbf{structured latent fingerprint} of artificial cognition. Just as molecular genetics enabled biology to transcend surface taxonomies and uncover causal mechanisms, we contend that a \textit{genomic lens} is now essential for machine learning--one that can \textbf{quantify}:


\begin{tcolorbox}
\begin{itemize}
    \item[\ding{93}] \textbf{Layer Importance and Semantic Specialization}: Not all layers contribute equally to a model’s epistemic structure. A growing body of evidence~\citep{belrose2023mechanistic, geva2022transformer, dai2023knowledge, liu2023hidden} reveals that semantic representations, cultural memory, and alignment behavior disproportionately concentrate in the mid-to-upper transformer layers--particularly the final 10 layers in $~30$-layer models. These layers encode more than surface patterns; they carry deep \emph{semantic priors} and value shifts induced by alignment, fine-tuning, and cultural adaptation. For \textbf{\textit{n}}{\fontfamily{ptm}\selectfont DNA} to serve as a meaningful genomic diagnostic, it must trace inheritance, drift, and trait transformation across these epistemically sensitive regions.

    \item[\ding{93}] \textbf{Semantic Drift and Heritable Traits}: Subtle misalignments and persistent divergences--documented in alignment studies~\citep{zhou2023alignmentdrift, ganguli2023reducing}--can occur even when models appear behaviorally consistent. These are not superficial perturbations but inheritable epistemic traits passed along neural offspring~\citep{wu2024seamless, xu2023aligning}.

    \item[\ding{93}] \textbf{Value Simulation vs. Internalization}: As models grow more context-sensitive, they learn to \emph{simulate} alignment without embodying its values~\citep{jacobs2024evalaware, perez2022discovering}. Disentangling true normative internalization from strategic mimicry is essential for any meaningful epistemic inspection.

    \item[\ding{93}] \textbf{Plasticity and Collapse}: Aggressive fine-tuning, distillation, or ideological merging can induce \emph{plasticity collapse}--a reduction in epistemic flexibility and semantic richness~\citep{liu2023lost, bai2023constitutional}. This demands metrics that trace both robustness and degeneration over time.

    \item[\ding{93}]\textbf{Latent Cultural Conflict}: In multilingual or cross-cultural settings, models often encode conflicting or incoherent value systems~\citep{mukherjee2024inconsistency, chen2023helpfulness}. These conflicts are not visible through BLEU or ROUGE--they reside in the model’s latent belief structure and must be surfaced through geometric lineage analysis.

    \item[\ding{93}] \textbf{Topological Continuity}: Alignment and fine-tuning warp the internal geometry of models in nontrivial ways~\citep{chiang2023can, liu2023hidden}. \textbf{\textit{n}}{\fontfamily{ptm}\selectfont DNA} must preserve continuity and interpretability of trajectories across such transformations.

    \item[\ding{93}] \textbf{Epistemic Mutation}: Merging preferences, annotator distributions, or learned behaviors--as explored by~\cite{bakker2024uniting}--creates emergent traits that standard metrics cannot track. These mutations are only diagnosable through a genomic lens on representation evolution.
\end{itemize}
\end{tcolorbox}

\vspace{2mm}
\textbf{\textit{n}}\textnormal{DNA} empowers us to interrogate the \emph{hidden geometry} of learning--revealing how foundational operations such as \textbf{alignment}, \textbf{fine-tuning}, \textbf{quantization}, \textbf{pruning}, and \textbf{multilingual fusion} subtly but systematically reshape a model’s \emph{semantic core}. It uncovers \textbf{cultural instabilities} introduced through regional adaptation, traces \textbf{asymmetric inheritance} patterns across neural offspring, visualizes \textbf{latent reorganizations} induced by merging or distillation, and quantifies a model’s capacity to \emph{resist} or \emph{absorb} conflicting epistemic pressures.


These phenomena--often dismissed as quirks--are in fact \emph{heritable traits}, etched into the model’s internal manifold. When viewed through this lens, \emph{model collapse}, \emph{alignment-induced drift}, and \emph{semantic mimicry} cease to be incidental failures and instead emerge as structural signatures of deeper latent dynamics. \textbf{\textit{n}}\textnormal{DNA} thus transcends metaphor to become a \textbf{scientific grammar} for measuring \emph{epistemic resilience}, \emph{semantic coherence}, \emph{cultural consistency}, and \emph{trait inheritance}--offering a principled lens through which to \textbf{govern}, \textbf{understand}, and \textbf{audit} the evolving anatomy of artificial cognition.



We further posit that \textbf{cultural provenance} induces a distinct \emph{layerwise calibration effect}, predominantly localized in the final decoder layers $\ell \in [20, 30]$, where sociolinguistic priors exert the strongest influence on output distribution. To capture this, we introduce the \textbf{nDNA Score}--a composite diagnostic unifying: \textbf{(i)} \emph{Spectral curvature} $\kappa_\ell$, reflecting the compression and warping of conceptual flow; \textbf{(ii)} \emph{thermodynamic length} $\mathcal{L}_\ell$, quantifying the epistemic effort required to traverse belief transitions; and \textbf{(iii)} the norm of the \emph{Belief Vector Field} $\|\mathbf{v}_\ell^{(c)}\|$, measuring the directional intensity of latent cultural drift.


Together, these dimensions form a latent semantic fingerprint--a high-dimensional, biologically inspired signature of internal cognition--enabling us to \textbf{trace}, \textbf{compare}, and \textbf{govern} the \emph{neural evolution} of foundation models with unprecedented granularity.





\subsection{Spectral Curvature ($\kappa_\ell$): A Geometric Lens on Latent Bending}

\noindent
\textbf{What is spectral curvature?} In classical geometry, curvature quantifies how much a path deviates from being straight--measuring local bending of a trajectory. In \textbf{spectral geometry} and \textbf{harmonic analysis}, curvature extends to how signals or paths behave in frequency space or under operators that encode structure (e.g., Laplacians, difference operators). \emph{Spectral curvature} refers to curvature derived through such operators--capturing the \emph{shape of latent signals} as they evolve across layers of a model.

\vspace{0.8em}

\noindent
\textbf{Why spectral for latent manifolds?}  
In foundation models, hidden representations form a sequence of activations $\{ h_\ell \}_{\ell=0}^L$ across layers. These representations trace a path in high-dimensional latent space. The \emph{shape} of this path encodes the model’s \textbf{internal conceptual flow}--how its beliefs evolve as it integrates priors, inputs, and alignment constraints. \textbf{Spectral operators} (such as discrete Laplacians or difference operators) naturally quantify how this path bends or accelerates--making them ideal for probing internal geometry. Unlike mere distance measures, \emph{spectral curvature} reflects \textbf{intrinsic shape}, invariant under reparameterization.

\vspace{0.8em}

\noindent
\textbf{Formulation and derivation.}  
Consider hidden activations $h_\ell \in \mathbb{R}^d$ at each layer $\ell$. The \textbf{first-order difference}
\[
\Delta h_\ell := h_\ell - h_{\ell-1}
\]
approximates the local directional change of latent states--a discrete analogue of \emph{velocity} in latent space.

\vspace{0.5em}

\noindent
To capture bending, we compute the change in this directional flow--the \textbf{second-order difference}:
\[
\Delta^2 h_\ell := \Delta h_{\ell+1} - \Delta h_\ell = (h_{\ell+1} - h_\ell) - (h_\ell - h_{\ell-1}) = h_{\ell+1} - 2 h_\ell + h_{\ell-1}.
\]
This operator acts like a \emph{discrete Laplacian} along the latent path, highlighting where the model’s \textbf{internal belief flow} deviates from a straight trajectory.

\vspace{0.5em}

\begin{tcolorbox}[
  colback=gray!5!white,
  colframe=black!75!white,
  boxrule=0.5pt,
  arc=2pt
]
\centering
\emph{Spectral curvature at layer $\ell$ is defined as:}
\[
\kappa_\ell := \big\| \Delta^2 h_\ell \big\| = \big\| h_{\ell+1} - 2 h_\ell + h_{\ell-1} \big\|
\]
\end{tcolorbox}

\vspace{0.5em}

\noindent
In continuous form, this corresponds to:
\[
\kappa(s) = \left\| \frac{d^2 h(s)}{ds^2} \right\|
\]
where $s$ parameterizes depth through the network. Our discrete $\kappa_\ell$ provides a practical, layerwise estimator.

\vspace{0.8em}


\noindent
\textbf{Why is this meaningful?}  
Peaks in $\kappa_\ell$ mark layers where internal geometry is most dynamic--zones of \emph{semantic inflection}, \emph{belief compression}, or \emph{ideological absorption}. These are the structural signatures of \textbf{internal epistemic adaptation}, essential to trace cultural inheritance and alignment drift.

\vspace{0.3em}

\noindent
\textbf{Lineage and context.}  
Spectral curvature builds on tools from \textbf{geometric deep learning}, \textbf{equivariant architectures}, \textbf{Ricci flow in machine learning}, and \textbf{spectral graph analysis}~\citep{farzam2024ricci, cho2023mixedcurvature, gasteiger2021gemnet, xu2022spherical, konf2021hierarchical, ying2021equivariant, hu2022lie, hess2023spectral, wang2021geomtransformer, raposo2023spectral}. Within \textbf{\textit{n}}\textnormal{DNA}, it serves as a \textbf{principled geometric fingerprint}--revealing not only \emph{what} is encoded, but \emph{how} internal belief pathways are reshaped to encode it. Figure~\ref{fig:spectral_curvature} illustrates how spectral curvature \(\kappa_\ell\) measures second-order geometric changes in latent representations across transformer layers, revealing critical semantic inflection points that reflect nuanced, layer-specific restructuring in belief and ideologically influenced model epistemic.


\begin{figure}[H]
    \centering
    \includegraphics[width=\linewidth]{Figures/spectral.png}
    \includegraphics[width=\linewidth]{Figures/spectral_3d.png}
    \vspace{-4mm}
    \caption{\textbf{Spectral Curvature} ($\boldsymbol{\kappa_\ell}$) quantifies second-order deviations in latent representations across transformer layers--computed via the discrete geometric operator $\boxed{\kappa_\ell := \| h_{\ell+1} - 2 h_\ell + h_{\ell-1} \|}$. High curvature signals \emph{semantic inflection points} where internal geometry bends sharply--often in \textbf{culturally dense}, \textbf{ideologically loaded}, or \textbf{epistemically volatile} regions. Peaks in $\kappa_\ell$ typically emerge in upper decoder layers ($\ell \in [21,30]$), where the model accommodates sociolinguistic priors during alignment, multicultural or multilingual fusion. Within the \textbf{\textit{n}}\textnormal{DNA} framework, such curvature reflects \emph{latent inheritance dynamics}, offering a fine-grained geometric fingerprint of representational restructuring.}
    \label{fig:spectral_curvature}
\end{figure}



\subsection{Thermodynamic Length ($\mathcal{L}_\ell$): Epistemic Effort Across Layers}

\noindent
\textbf{What is thermodynamic length?}  
In \textbf{statistical thermodynamics} and \textbf{information geometry}, \emph{thermodynamic length} measures the cumulative effort--or “\emph{work}”--required for a system to transition between states on a statistical manifold. It integrates local gradient energy along a trajectory, providing an \emph{intrinsic cost measure} that is independent of parametrization.

\vspace{0.5em}

\noindent
\textbf{Why thermodynamic length for foundation models?}  
In foundation models, layers trace a \textbf{path through latent belief space}. 
As input data and alignment priors reshape activations, the model expends internal \textbf{computational effort} to adjust its belief state. 
\emph{Thermodynamic length quantifies this latent effort} --- measuring not just \emph{what} the model knows, but \emph{how hard} it works to adapt that knowledge across layers in response to epistemic pressures (e.g., cultural fusion, alignment shifts).

\vspace{0.3em}

\noindent
\textbf{Mathematical intuition.}  
Let $h_\ell$ denote the latent state at layer $\ell$, and $\mathcal{M}$ the model’s latent manifold. 
Layer transitions define a curve 
$\gamma: [0,L] \to \mathcal{M}$ 
whose thermodynamic length is
\[
\boxed{
\mathcal{L}(\gamma) = \int_0^L \sqrt{ \big\langle \dot{\gamma}(s),\, \mathcal{G}_{\mathrm{Fisher}}\, \dot{\gamma}(s) \big\rangle }\, ds
}
\]
where $\mathcal{G}_{\mathrm{Fisher}}$ is the Fisher information metric. 
Here, $\mathcal{L}(\gamma)$ represents the \emph{intrinsic work} needed to traverse $\gamma$ on $\mathcal{M}$.

\vspace{0.3em}

\noindent
\textbf{Interpretation.}  
High thermodynamic length indicates regions where latent geometry \textbf{stretches} --- where the model’s belief space undergoes substantial reconfiguration to reconcile priors and input. 
This formalism reveals \emph{not just where} latent states change, but the \emph{cost structure of that change}. 
Zones of large $\mathcal{L}_\ell$ mark points of \textbf{alignment tension}, \textbf{cultural fusion}, or \textbf{complex reasoning}, where internal scaffolds are under maximum stress.

\vspace{0.3em}

\noindent
\emph{Thermodynamic length offers a window onto the model’s ``latent energy budget'' --- illuminating how internal belief states reshape to meet complexity, constraint, and context.}


\vspace{0.8em}

\noindent
\textbf{Formulation.}  
Let $p_\ell(y|x)$ denote the model’s conditional distribution at layer $\ell$ given input $x$. 
The local epistemic cost is reflected in the squared norm of the gradient of log-likelihood with respect to model parameters:
\[
\big\| \nabla_\theta \log p_\ell(x) \big\|^2.
\]
This quantity measures how much the model must \emph{adjust its parameters locally} at layer $\ell$ to improve its fit to input $x$. \emph{Thermodynamic length at layer $\ell$} aggregates this cost across the dataset $\mathcal{D}$:
\begin{tcolorbox}[
  colback=gray!5!white,
  colframe=black!75!white,
  boxrule=0.5pt,
  arc=2pt
]
\centering
\emph{Thermodynamic length at layer $\ell$ is defined as:}
\[
\mathcal{L}_\ell := \sum_{x \in \mathcal{D}} \big\| \nabla_\theta \log p_\ell(x) \big\|^2
= |\mathcal{D}| \, \mathbb{E}_{x \sim \mathcal{D}} \big\| \nabla_\theta \log p_\ell(x) \big\|^2.
\]
\end{tcolorbox}

\noindent
This formulation reveals that $\mathcal{L}_\ell$ captures both the \emph{average local effort} and its scaling with dataset size.  
Furthermore, in differential geometric terms, thermodynamic length can be written as a path energy:
\[
\mathcal{L}_\ell = \int_{\gamma_\ell} 
\left\langle \frac{d h_\ell}{ds}, 
\mathcal{G}_{\mathrm{Fisher}}(h_\ell)
\frac{d h_\ell}{ds} \right\rangle ds
\]
where $h_\ell$ denotes latent trajectories at layer $\ell$, $\mathcal{G}_{\mathrm{Fisher}}$ the Fisher information metric, and $s$ arc length along $\gamma_\ell$. Thus, $\mathcal{L}_\ell$ can be seen as an \emph{energy integral over the belief manifold} -- capturing how much internal ``\emph{heat}'' or computational work is generated to reconcile prior belief state with new input at depth $\ell$.


\noindent
\textbf{Why is this meaningful?}  
Unlike static capacity metrics or weight magnitudes, $\mathcal{L}_\ell$ is \emph{dynamically grounded}: it measures where the model actively strains to reconcile competing epistemic demands. In regions of high $\mathcal{L}_\ell$, the model’s \textbf{latent geometry} is under tension--\emph{reshaping itself} to accommodate alignment constraints, cultural priors, or multilingual semantics.

\vspace{0.3em}

\noindent
\textbf{Lineage and context.}  
This diagnostic builds on the \textbf{Fisher–Rao metric} in \textbf{information geometry} and \textbf{thermodynamic length formalism} from statistical physics~\citep{crooks2007measuring, oliviero2023thermodynamics, farzam2024ricci, wagner2023thermodynamic}. Thus \textbf{\textit{n}}\textnormal{DNA} provides a \emph{complementary view} to spectral curvature--capturing not where the model bends, but \emph{how hard it works} to do so. Together, these axes form a \textbf{neurogeometric anatomy} of latent belief adaptation. 

Figure~\ref{fig:thermo_length} shows thermodynamic length \(\mathcal{L}_\ell\), quantifying latent epistemic effort and semantic restructuring across transformer layers.



\begin{figure}[H]
    \centering
    \includegraphics[width=\linewidth]{Figures/thermodynamics.png}
    \includegraphics[width=0.9\linewidth]{Figures/thermodynamics_3d.png}
    \vspace{-6mm}
    \caption{\textbf{Thermodynamic Length} $\boxed{\mathcal{L}_\ell := \sum_{x \in \mathcal{D}} \big\| \nabla_\theta \log p_\ell(x) \big\|^2}$ quantifies the \emph{epistemic work} performed across transformer layers, calculated as the \textbf{cumulative squared gradient norm of layerwise log-likelihoods}. Higher values signal \emph{internal resistance}--zones of significant restructuring, belief compression, or negotiation of conflicting priors. In culturally fine-tuned models, these peaks localize to upper decoder layers, indicating intense adaptation near output-generating blocks. Within the \textbf{\textit{n}}\textnormal{DNA} construct, $\mathcal{L}_\ell$ helps reveal latent epistemic effort that underlies surface-level behavior. This metric thus provides a nuanced window into where and how models internally allocate effort during learning and inference.}
    \label{fig:thermo_length}
\end{figure}

\vspace{0.5em}




\subsection{Belief Vector Field ($\mathbf{v}_\ell^{(c)}$): Cultural Drift in Latent Space}

\noindent
\textbf{What is the Belief Vector Field} --  
In \textbf{differential geometry} and \textbf{physics}, a \emph{vector field} describes a directional force applied at each point of a space. Inspired by this, the \textbf{Belief Vector Field} models the \emph{directional semantic force} that a specific culture or value system exerts on a model’s latent representations. It encodes \emph{where}, \emph{how strongly}, and \emph{in what direction} cultural priors act within the model’s internal geometry--functioning as a \textbf{semantic compass} through the latent manifold.

\vspace{0.5em}

\noindent
\textbf{Why a vector field for cultural influence?}  
While \textbf{spectral curvature} ($\kappa_\ell$) captures how sharply latent paths bend, and \textbf{thermodynamic length} ($\mathcal{L}_\ell$) how hard the model works during adaptation, neither tells us the \emph{source}, \emph{direction}, or \emph{origin} of that adaptation. The Belief Vector Field offers this missing piece: it traces the latent steering aka torison applied by culture-conditioned priors--\emph{where the model is being pushed in latent space, by what epistemic force, and toward which semantic direction}. This makes it a critical diagnostic for studying \textbf{cultural drift}, \textbf{ideological imprinting}, and \textbf{alignment tension}.





\begin{figure}[H]
    \centering
    \includegraphics[width=\linewidth]{Figures/belief_vector_field_healthy_static_annotated.png}
    \vspace{-8mm}
    \caption{
    \textbf{Belief Vector Field Visualization}: 
    $\mathbf{v}_\ell^{(c)} = 
    \mathbb{E}_{x \sim \mathcal{P}^{(c)}_{\mathrm{CIVIC}}}
    \left[
    \nabla_{h_\ell} \log p(y \mid x)
    \right]$
    represents the \emph{belief semantic steering force} at layer~$\ell$ toward concept~$c$, conditioned on CIVIC cultural priors (cf.~\cref{sec:aether_benchmark}).
    \textbf{Large magnitudes} (e.g., $\| \mathbf{v}_\ell^{(c)} \| \in [0.15, 0.50]$) indicate \emph{strong directional pressure}--zones where cultural values actively reshape latent geometry.
    \textit{Color-coded arrows} trace distinct conceptual trajectories (\textcolor[rgb]{0.29,0.00,0.51}{protest}, \textcolor[rgb]{0.00,0.00,0.50}{peace}, \textcolor[rgb]{0.00,0.50,0.50}{order}, \textcolor[rgb]{0.24,0.70,0.44}{power}, \textcolor[rgb]{0.42,0.56,0.14}{disobedience}, \textcolor[rgb]{0.85,0.65,0.13}{justice}), while numeric labels quantify local steering strength.
    Upper layers ($\ell \ge 20$) typically exhibit \textbf{epistemic reorientation}, where cultural priors most heavily influence belief encoding.
    Such visualizations reveal whether a model internalizes culturally contingent reasoning or merely mimics alignment at the output surface.
    }
    \label{fig:belief_vector_field}
\end{figure}

\noindent
\textbf{Formulation and derivation.}  
Let $p(y|x)$ denote the model’s conditional output distribution for input $x$, and let $h_\ell$ be the latent representation at layer $\ell$. The local belief gradient, $\nabla_{h_\ell} \log p(y|x)$, measures how a small change in $h_\ell$ would affect output confidence--a proxy for \emph{semantic force} at that layer. To extract the culturally conditioned semantic force, we compute its expectation over a culture-specific distribution $\mathcal{P}^{(c)}$:
\begin{tcolorbox}[
  colback=gray!5!white,
  colframe=black!75!white,
  boxrule=0.5pt,
  arc=2pt
]
\centering
\emph{Belief vector field at layer $\ell$ for a given manifold condition is defined as:}
\[
\mathbf{v}_\ell^{(c)} := \mathbb{E}_{x \sim \mathcal{P}^{(c)}} \left[ \nabla_{h_\ell} \log p(y|x) \right]
\]
\end{tcolorbox}
where $\mathcal{P}^{(c)}$ represents inputs emblematic of givem manifold condition $c$ (e.g., regional, linguistic, ideological contexts). This formulation captures not just latent deformation, but \emph{its cause}: how cultural priors exert directional influence within the belief manifold.

\noindent
\textbf{Why is this meaningful?}  
$\mathbf{v}_\ell^{(c)}$ provides a directional lens on latent dynamics. High $\| \mathbf{v}_\ell^{(c)} \|$ signals regions where the model is \emph{actively redirected} by external cultural forces--offering diagnostic power for detecting \textbf{ideological drift}, \textbf{semantic conflict}, or \textbf{bias inheritance}. Unlike $\kappa_\ell$ or $\mathcal{L}_\ell$, which capture internal geometry, $\mathbf{v}_\ell^{(c)}$ reveals \emph{external epistemic pressure} and its directional impact.

\vspace{0.3em}

\noindent
\textbf{Lineage and context.}  
This diagnostic builds upon belief geometry, alignment drift studies, and cultural bias tracing in NLP~\citep{wang2023culturalbias, zhou2023alignmentdrift, shen2023beliefgeometry, arora2023stereoset, bommasani2023foundation, peng2024cultural, laurens2024anthropic, kang2024biasfairness, de2023latentbias, gao2023value}. Within the \textbf{\textit{n}}\textnormal{DNA} construct, it integrates with curvature and length to offer a holistic neurogeometric portrait--revealing \emph{how}, \emph{why}, and \emph{where} foundation models inherit, adapt, or distort beliefs under cultural influence.

\vspace{0.3em}

\noindent
\textbf{Interpretability in practice.}  
By mapping $\mathbf{v}_\ell^{(c)}$ across layers and cultures, we can trace \textbf{cultural provenance}, identify \textbf{ideological pressure zones}, and diagnose \textbf{inheritance asymmetry} in multilingual or aligned models. This directional fingerprint informs audits of model bias, robustness, and alignment integrity--providing the missing vectorial dimension in understanding machine cognition.


\subsection{\textbf{\textit{n}}\textnormal{DNA}: Unified Epistemic Inheritance Measure}

\noindent
\textbf{Why a unified score?}  
While \textbf{spectral curvature} ($\kappa_\ell$), \textbf{thermodynamic length} ($\mathcal{L}_\ell$), and the \textbf{belief vector field norm} ($\| \mathbf{v}_\ell^{(c)} \|$) each offer unique insight into latent dynamics, they operate on distinct facets of \emph{epistemic geometry}:

\vspace{0.8em}

\begin{tcolorbox}[
  colback=blue!3!white,
  colframe=black!80!white,
  boxrule=0.6pt,
  arc=3pt,
  width=\textwidth
]
\small
\textbf{The nDNA score is a cumulative measure of latent geometry, quantifying how a large language model adapts its internal scaffolding to a given corpus.} 
It integrates three key components at each layer $\ell$:
\begin{itemize}[leftmargin=1.5em]
  \item \textbf{Curvature} ($\kappa_\ell$): how \emph{twisted} or \emph{bent} the latent manifold is; captures \emph{how sharply} internal trajectories bend -- a scalar measure of latent acceleration.
  \item \textbf{Length} ($\mathcal{L}_\ell$): how much \emph{latent work} or displacement occurs as representations evolve; quantifies \emph{how hard} the model works to adapt its beliefs -- a scalar effort integral.
  \item \textbf{Belief vector norm} ($\|\mathbf{v}_\ell^{(c)}\|$): how \emph{strong} the model’s belief signal is for that corpus; encodes \emph{where} and \emph{how strongly} cultural priors steer latent space -- a scalar magnitude derived from the vector field.
\end{itemize}
\begin{center}
\emph{Formally, we define the \textbf{\textit{n}}\textnormal{DNA} score as:}
\[\boxed{
\texttt{nDNA} := \sum_{\ell=1}^{L} \omega_\ell \cdot \kappa_\ell \cdot \mathcal{L}_\ell \cdot \| \mathbf{v}_\ell^{(c)} \|}
\]
\end{center}
\end{tcolorbox}



\vspace{0.5em}

\begin{comment}
\begin{itemize}[leftmargin=1.5em]
    \item $\kappa_\ell$ captures \emph{how sharply} internal trajectories bend -- a scalar measure of latent acceleration.
    \item $\mathcal{L}_\ell$ quantifies \emph{how hard} the model works to adapt its beliefs -- a scalar effort integral.
    \item $\| \mathbf{v}_\ell^{(c)} \|$ encodes \emph{where} and \emph{how strongly} cultural priors steer latent space -- a scalar magnitude derived from the vector field.
\end{itemize}
\end{comment}


Individually, these measures illuminate latent strain, adaptation cost, and cultural pressure. But to assess \emph{inheritance as a whole} -- how traits propagate through \textbf{fine-tuning}, \textbf{merging}, or \textbf{distillation} -- we must integrate these into a single diagnostic that reflects combined latent geometry, epistemic work, and directional influence.

\vspace{0.8em}

\noindent
\textbf{Designing the composite measure.}  
Since $\kappa_\ell$ and $\mathcal{L}_\ell$ are scalars, and $\| \mathbf{v}_\ell^{(c)} \|$ reduces directional drift to scalar magnitude, their product forms a natural joint measure of:
\emph{internal bending} ($\kappa_\ell$),
\emph{internal epistemic effort} ($\mathcal{L}_\ell$),
and \emph{external drift pressure} ($\| \mathbf{v}_\ell^{(c)} \|$).
To balance their contributions across depth, we introduce layer weights $\omega_\ell$, emphasizing semantically active or epistemically significant layers (e.g., $\omega_\ell$ higher in upper decoder blocks).

\clearpage
\newpage
\begin{figure*}[ht!]
\centering
\includegraphics[width=0.98\textwidth]{Figures/ndna_refined_story_finalframe.png}
\vspace{-4mm}
\caption{
\textbf{The compositional anatomy of neural DNA (\emph{n}DNA) through curvature, length, and belief geometry.}
This figure illustrates how \emph{n}DNA arises as a layered product of three latent quantities. 
First, \textbf{spectral curvature} $\boldsymbol{\kappa_\ell}$ measures latent manifold bending and flexibility (latent acceleration), indicating how sharply the internal geometry twists at layer $\ell$. 
Second, \textbf{thermodynamic length} $\boldsymbol{\mathcal{L}_\ell}$ quantifies the accumulated epistemic effort (latent adaptation energy) and reflects how hard the model works to reconcile prior beliefs with new input and alignment signals. 
Third, \textbf{belief vector norm} $\|\mathbf{v}_\ell^{(c)}\|$ encodes the magnitude of latent directional force imposed by corpus priors or alignment signals. 
The joint trajectory in $(\kappa_\ell, \mathcal{L}_\ell, \|\mathbf{v}_\ell^{(c)}\|)$ space, color-coded by the composite score, shows how bending, effort, and steering co-evolve across layers. 
The combined latent signature is formalized as $\mathit{nDNA}_\ell = \kappa_\ell \cdot \mathcal{L}_\ell \cdot \|\mathbf{v}_\ell^{(c)}\| = 0.0024$ (example layer), with high values identifying zones of intense latent reconfiguration where geometry and adaptation forces align. 
Color-keyed descriptors (``Latent bending'', ``Epistemic effort'', ``Belief steering'') guide visual interpretation. 
The figure illustrates how large language models coordinate latent bending, effort, and steering to build a neurogeometric scaffold that adapts flexibly to task complexity while remaining anchored in a universal latent structure.
}
\label{fig:ndna_story}
\end{figure*}
%\clearpage
%\newpage





\noindent
This composite score integrates scalar and vector-derived diagnostics into a unified measure of \emph{epistemic inheritance} -- quantifying the latent structure and cultural traits a model carries forward from its neural ancestry.

\vspace{0.3em}

\noindent
\textbf{Rationale for multiplicative integration.}  
This form spotlights layers where latent paths bend sharply, belief adaptation incurs significant effort, and cultural or alignment pressures apply strong directional force. 
High scores identify zones of \emph{intense latent reconfiguration}, where internal dynamics and external pressures converge to reshape the model’s reasoning space.


\vspace{0.3em}

\noindent
\textbf{Role of $\omega_\ell$.}  
The weight $\omega_\ell$ serves as a lens to prioritize semantically expressive, epistemically active regions of the network. It may be set uniformly, hand-tuned, or optimized against alignment drift benchmarks, bias metrics, or interpretability objectives.

\vspace{0.3em}

\noindent
\textbf{Interpretability and utility.}  
The $\texttt{nDNA}$ score provides a compact fingerprint of model inheritance:
\begin{itemize}[leftmargin=1.5em]
    \item It enables direct comparison of parent and child models post \textbf{fine-tuning}, \textbf{merging}, or \textbf{distillation}.
    \item It highlights zones of \emph{semantic mutation}, \emph{ideological absorption}, or \emph{cultural drift}.
    \item It serves as a proxy for \emph{latent epistemic integrity} -- quantifying the hidden cost and directionality of neural evolution.
\end{itemize}

\vspace{0.3em}

\noindent
\textbf{Conviction.}  
By unifying \textbf{spectral}, \textbf{thermodynamic}, and \textbf{vectorial} diagnostics, the $\texttt{nDNA}$ score functions as a \textbf{heritable geometry index} -- diagnosing how latent traits persist, mutate, or degrade as foundation models evolve.



\subsection{nDNA Geometry - A closer Look}

\noindent
The notion of \textbf{nDNA} arises from a simple yet profound insight: modern foundation models do not merely produce outputs--they embody a latent cognitive structure that governs how they reason, adapt, and evolve~\citep{bommasani2023foundation, ganguli2023reducing}. This latent structure is not directly encoded in model weights or activations alone; rather, it emerges in the internal geometry of belief formation, semantic flow, and epistemic adaptation across layers~\citep{liu2023hidden, wang2021geomtransformer}. We define the \textbf{nDNA geometry} of a model as the joint distribution of its 
\textbf{spectral curvature} ($\boldsymbol{\kappa_\ell}$), 
\textbf{thermodynamic length} ($\boldsymbol{\mathcal{L}_\ell}$), and 
\textbf{belief vector field norm} ($\| \mathbf{v}_\ell^{(c)} \|$)
layer-by-layer. 
This triad forms a high-dimensional semantic fingerprint that encodes a model's \emph{inheritance stability}, 
\emph{alignment dynamics}, 
and \emph{cultural drift}---analogous to how biological DNA records heritable traits and mutations~\citep{shen2023beliefgeometry, bakker2024uniting}.









\begin{table}[H]
\centering
\caption{
An \textbf{illustrative nDNA example}: that captures the \emph{semantic genome} of a foundation model through the joint interplay of \textbf{spectral curvature} ($\boldsymbol{\kappa_\ell}$), \textbf{thermodynamic length} ($\boldsymbol{\mathcal{L}_\ell}$), \textbf{belief vector norm} ($\| \mathbf{v}_\ell^{(c)} \|$) across layers. Each of these quantities offers a distinct geometric and epistemic lens: $\boldsymbol{\kappa_\ell}$ measures the \emph{local acceleration} of latent representations, $\boldsymbol{\mathcal{L}_\ell}$ quantifies the cumulative \emph{internal work} required to traverse the belief manifold, while $\| \mathbf{v}_\ell^{(c)} \|$ encodes the \emph{magnitude of cultural drift} imposed on latent activations. The \emph{color intensities} shown alongside each value reflect relative magnitude within column-specific ranges: \textcolor[rgb]{0,0.6,0}{\rule{1em}{1em}} low, 
\textcolor[rgb]{0.8,0.8,0}{\rule{1em}{1em}} moderate, \textcolor[rgb]{1,0.5,0}{\rule{1em}{1em}} high, \textcolor[rgb]{0.8,0,0}{\rule{1em}{1em}} very high. For this example, spectral curvature spans $\boldsymbol{\kappa_\ell} \in [0.0400, 0.0700]$, thermodynamic length $\boldsymbol{\mathcal{L}_\ell} \in [0.80, 1.20]$, and belief vector norm $\| \mathbf{v}_\ell^{(c)} \| \in [0.55, 0.75]$--revealing regions where the \emph{latent manifold bends}, \emph{epistemic energy intensifies}, or \emph{external priors steer internal cognition}. This triad forms what we term the model's \textbf{nDNA}: a compact, high-dimensional \emph{semantic fingerprint} 
that encodes the hidden geometry of belief. 
It enables us to diagnose zones of \emph{inheritance stability}, detect \emph{ideological absorption}, and trace \emph{latent mutations} introduced by fine-tuning, alignment, or architectural choice. The pattern of these quantities across layers constitutes a signature as unique as a biological genome -- a map of how artificial cognition evolves, 
remembers, and adapts.
}
\begin{tabular}{|c|c|c|c|l|}
\hline
\textbf{Layer} 
& $\boldsymbol{\kappa_\ell}$ 
& $\boldsymbol{\mathcal{L}_\ell}$ 
& $\|\mathbf{v}_\ell^{(c)}\|$ 
& \textbf{Belief Vector} $\mathbf{v}_\ell^{(c)}$ \\
\hline
20 & \cellcolor{green!20}0.0412 & \cellcolor{yellow!30}0.9123 & \cellcolor{orange!20}0.6521 & $[0.1204, -0.0502, 0.0896, \ldots, 0.0402]$ \\
21 & \cellcolor{green!30}0.0458 & \cellcolor{green!15}0.8123 & \cellcolor{red!30}0.7523 & $[0.1301, -0.0351, 0.0950, \ldots, 0.0431]$ \\
22 & \cellcolor{yellow!20}0.0523 & \cellcolor{orange!30}1.0120 & \cellcolor{green!20}0.5823 & $[0.1423, -0.0312, 0.0994, \ldots, 0.0488]$ \\
23 & \cellcolor{orange!20}0.0581 & \cellcolor{yellow!20}0.9021 & \cellcolor{yellow!30}0.6912 & $[0.1534, 0.0270, 0.1042, \ldots, 0.0512]$ \\
24 & \cellcolor{red!20}0.0639 & \cellcolor{red!40}1.1023 & \cellcolor{green!15}0.5520 & $[0.1667, 0.0205, 0.1105, \ldots, 0.0543]$ \\
25 & \cellcolor{yellow!30}0.0505 & \cellcolor{orange!20}0.9420 & \cellcolor{red!40}0.8124 & $[0.1602, -0.0251, 0.1081, \ldots, 0.0504]$ \\
26 & \cellcolor{green!15}0.0398 & \cellcolor{green!20}0.8520 & \cellcolor{yellow!20}0.6120 & $[0.1251, 0.0450, 0.0912, \ldots, 0.0418]$ \\
27 & \cellcolor{yellow!20}0.0512 & \cellcolor{red!30}1.0520 & \cellcolor{orange!30}0.7222 & $[0.1455, -0.0322, 0.1005, \ldots, 0.0477]$ \\
28 & \cellcolor{orange!30}0.0590 & \cellcolor{yellow!30}0.9320 & \cellcolor{green!30}0.5721 & $[0.1577, 0.0285, 0.1078, \ldots, 0.0499]$ \\
29 & \cellcolor{red!30}0.0672 & \cellcolor{orange!30}1.0123 & \cellcolor{yellow!20}0.6322 & $[0.1701, -0.0198, 0.1142, \ldots, 0.0533]$ \\
30 & \cellcolor{orange!20}0.0555 & \cellcolor{green!15}0.8221 & \cellcolor{red!20}0.7720 & $[0.1620, -0.0242, 0.1101, \ldots, 0.0510]$ \\
\hline
\end{tabular}
\label{tab:ndna_example}
\end{table}








\clearpage














\noindent
Table~\ref{tab:ndna_example} provides an \emph{illustrative example of nDNA geometry}, highlighting how these quantities vary across depth in a representative model. 
Rather than simple monotonic trends, we observe intricate layer-wise patterns: certain layers exhibit elevated curvature ($\kappa_\ell > 0.06$), signaling sharp latent reorientation~\citep{cho2023mixedcurvature}, while others concentrate thermodynamic length ($\mathcal{L}_\ell > 1.10$), reflecting zones of intense internal work to reconcile competing priors~\citep{crooks2007measuring, oliviero2023thermodynamics}. 
The belief vector norm $\| \mathbf{v}_\ell^{(c)} \|$ exposes the directional cultural force acting on the latent manifold~\citep{peng2024cultural, zhou2023alignmentdrift}, marking layers where external alignment or sociolinguistic conditioning exerts greatest influence. 
Together, these values form a geometry-specific trace that distinguishes models by their latent adaptation history.



\section{The Corpus Dependence of nDNA: A Necessary Feature, Not a Flaw}

In biological systems, DNA is celebrated as the \emph{universal code of life} -- a sequence of nucleotides that, across all known organisms, governs the development, function, and inheritance of traits \citep{alberts2014molecular, lewin2013genes}. 
Yet despite this \textbf{universal structure}, the functional expression of DNA is profoundly \textbf{context-dependent}. The same genome, when expressed in different cellular contexts, gives rise to vastly different phenotypes: for instance, \textit{neurons} and \textit{hepatocytes} arise from identical genetic material yet serve radically different functions \citep{bird2007perceptions, davidson2006gene}. 
This context-sensitive expression is orchestrated through layered regulatory mechanisms, including \textbf{epigenetic modifications} \citep{bird2007perceptions}, \textbf{transcription factor (TF) binding} \citep{lambert2018human}, and \textbf{chromatin architecture remodeling} \citep{clapier2017mechanisms, dekker2013exploring}. 
These mechanisms form a hierarchical, probabilistic regulatory network that determines gene expression patterns in response to developmental and environmental cues \citep{alon2006introduction}. Figure \ref{fig:enhanced_dna_ndna} illustrates a hierarchical regulatory framework where universal DNA undergoes epigenetic modifications and context-specific transcription factor actions to produce specialized gene expression programs. Analogously, in large language models, this layered structure parallels \textbf{\textit{n}}DNA latent scaffolding that encodes both \emph{universal priors} and \emph{task-dependent adaptations}, enabling \textbf{coherent, flexible, and robust functional diversity} across domains.



\begin{figure}[ht!]
\centering
\begin{tikzpicture}[node distance=1.2cm and 1.5cm, align=center, font=\small]

% Universal DNA helix at the top
\node (dna_label) at (0,0.7) {\textbf{Universal DNA (shared genome)}};
\draw[decorate, decoration={coil, aspect=0.4, amplitude=4pt, segment length=6pt}, thick, blue] (0,0.5) -- (0,-0.5);
\node (dna) at (0,-0.5) {};

% Epigenetic regulators
\node (methyl) [rectangle, draw, fill=gray!20] at (-3,-2) {DNA methylation \\ (suppression)};
\node (acetyl) [rectangle, draw, fill=green!20] at (3,-2) {Histone acetylation \\ (activation)};
\draw[->, thick, gray] (dna) -- (methyl);
\draw[->, thick, green!70!black] (dna) -- (acetyl);

% TF binding
\node (tf_neuron) [rectangle, draw, fill=green!10, below=1.5cm of methyl] {Neuron TFs \\ (NeuroD, REST)};
\node (tf_hep) [rectangle, draw, fill=orange!10, below=1.5cm of acetyl] {Hepatocyte TFs \\ (HNF4, C/EBP$\alpha$)};
\draw[->, thick] (methyl) -- (tf_neuron);
\draw[->, thick] (acetyl) -- (tf_hep);

% Chromatin states
\node (chrom_neuron) [rectangle, draw, dashed, below=1.2cm of tf_neuron] 
{Open chromatin \\ neuron genes on};
\node (chrom_hep) [rectangle, draw, dashed, below=1.2cm of tf_hep] 
{Compact chromatin \\ neuron genes off};
\draw[->] (tf_neuron) -- (chrom_neuron);
\draw[->] (tf_hep) -- (chrom_hep);

% Gene programs
\node (gene_neuron) [rectangle, draw, rounded corners=3pt, fill=green!5, below=1cm of chrom_neuron] 
{Synaptic genes ON};
\node (gene_hep) [rectangle, draw, rounded corners=3pt, fill=orange!5, below=1cm of chrom_hep] 
{Detox/metabolic genes ON};
\draw[->] (chrom_neuron) -- (gene_neuron);
\draw[->] (chrom_hep) -- (gene_hep);

% Functions
\node (func_neuron) [below=0.7cm of gene_neuron] {Function: signaling, plasticity};
\node (func_hep) [below=0.7cm of gene_hep] {Function: glucose metabolism, detox};

\end{tikzpicture}

\caption{
\textbf{A hierarchical view of universal DNA and context-sensitive gene expression, as a biological parallel to nDNA latent scaffolding in LLMs.}
This figure illustrates how the \emph{same genome} (depicted as a universal DNA helix at the top) produces distinct functional outcomes through a layered and structured regulatory architecture. 
The \textbf{first regulatory layer} consists of \textit{epigenetic modifications}, including DNA methylation (linked with gene silencing) and histone acetylation (linked with gene activation) 
\citep{bird2007perceptions, clapier2017mechanisms}. 
These modifications influence chromatin accessibility, setting the stage for context-specific transcriptional control.  
The \textbf{second layer} involves \textit{cell-type-specific transcription factors (TFs)} -- for example, NeuroD and REST in neurons, or HNF4 and C/EBP$\alpha$ in hepatocytes -- which bind regulatory DNA elements 
and integrate signaling cues to guide gene expression programs \citep{lambert2018human, davidson2006gene}.  
The \textbf{third layer} reflects the resultant chromatin state: open, transcriptionally permissive configurations in neurons for synaptic gene activation, versus compact, repressive configurations in hepatocytes where those genes are silent 
\citep{thurman2012accessible, dekker2013exploring}.  
Finally, this hierarchical regulatory control produces \textit{functionally specialized gene programs}: neurons activate synaptic plasticity and axon signaling genes; hepatocytes activate detoxification and glucose metabolism genes 
\citep{lewin2013genes, alon2006introduction}.  
This layered architecture provides a powerful biological analogy for \emph{nDNA in LLMs}. 
Just as DNA’s expression is shaped by regulatory logic rather than random variation, \textbf{nDNA encodes both universal priors (shared across tasks)} -- such as pretrained latent manifolds, attention mechanisms, and model architecture -- 
and \textbf{corpus-dependent latent scaffolding}, emerging as the model adapts to specific tasks or domains \citep{olah2020zoom, geva2021transformer, beltagy2020longformer}. 
The analogy emphasizes that corpus dependence in nDNA is not a weakness or artifact, but a reflection of meaningful task adaptation: 
\emph{structured variation grounded in universal latent geometry}. 
This scaffolding ensures LLMs achieve \textit{functional diversity} across tasks while maintaining \textbf{coherence, alignment, and generalization}, much like gene regulatory networks ensure appropriate cellular identity 
and function despite operating from a common genome blueprint \citep{alon2006introduction, davidson2006gene}.  
The figure highlights that both biological DNA and nDNA exhibit clarity through complexity: 
\textbf{layered, interpretable hierarchies enabling flexible, robust expression across contexts}.
}
\label{fig:enhanced_dna_ndna}
\vspace{-3mm}
\end{figure}



Similarly, in large foundation models, the \emph{neural DNA (nDNA)} -- a composite measure of latent geometry encompassing \textbf{spectral curvature} ($\kappa$) \citep{belkin2019reconciling}, 
\textbf{thermodynamic length} ($L$) \citep{still2012thermodynamic}, 
and \textbf{latent belief vector norms} \citep{olah2020zoom} -- exhibits both \textbf{universal structure} and \textbf{corpus-specific adaptation}. 
LLMs encode universal latent priors through pretraining: architectural invariances \citep{vaswani2017attention}, semantic manifolds \citep{mikolov2013distributed, bommasani2021opportunities}, and attention-based relational structures \citep{geva2021transformer}. 
However, when probed with different corpora -- such as mathematical reasoning benchmarks (e.g. GSM8K \citep{cobbe2021training}), dialogue datasets (e.g. MultiWOZ \citep{budzianowski2018multiwoz}), or encyclopedic QA (e.g. SQuAD \citep{rajpurkar2016squad}) -- 
the model activates distinct latent scaffolding, producing task-specific geometric pathways.



In both systems, \textbf{structured variation emerges as a necessity}: 
in \textbf{biology}, to produce \emph{functional diversity} across cell types; 
in \textbf{LLMs}, to scaffold \emph{reasoning} across tasks while maintaining \textbf{alignment} and \textbf{generalization} \citep{bommasani2021opportunities, cobbe2021training}. 
Like \textbf{tissue-specific gene expression}, \textbf{corpus-dependent nDNA scaffolding} follows precise, \emph{learned priors} rather than arbitrary variation. 
\textbf{Mathematical models} of both systems reduce to \emph{path integrals over conditional cost}: 
\[
\mathcal{S}(c) = \int_{\gamma_c} \mathcal{C}(h_\ell; c) ds
\]
where $\gamma_c$ is the pathway for \emph{context} $c$ (cell type or corpus), and $\mathcal{C}$ reflects \emph{regulatory} or \emph{loss cost}.

\vspace{-2mm}
\begin{quote}
\emph{\textbf{Where DNA differentiates cells}, nDNA differentiates reasoning. Both systems achieve \textbf{functional coherence} through context-dependent geometry anchored in \textbf{universal code}.}
\end{quote}

Despite their \emph{contextual variation}, both \textbf{DNA} and \textbf{nDNA} encode \textbf{universal structure} that stabilizes functional diversity. 
In \textbf{biology}, this universality is embodied in the \emph{genetic code}: the shared language of \textbf{codons}, \textbf{conserved regulatory motifs}, and \textbf{chromatin architectural principles} that ensure coherent development across tissues \citep{lewin2013genes, alberts2014molecular}. 
In \textbf{large language models}, nDNA’s universality arises from the \textbf{shared latent priors} learned during pretraining: \textbf{attention-based relational structures} \citep{vaswani2017attention}, \textbf{semantic manifolds} \citep{mikolov2013distributed}, and \textbf{transformer-invariant latent symmetries} \citep{bommasani2021opportunities}. 
These priors act as the \emph{``genomic grammar''} that binds task-specific latent pathways into a \textbf{coherent reasoning framework}.

\[
\boxed{
\textbf{DNA: } \Sigma^3 / \ker \phi \to \mathcal{A} \quad 
\textbf{nDNA: } \mathcal{X} / G_{\mathrm{LLM}} \to V/G
}
\]

Such \textbf{universal structure} enables \textbf{generalization}: 
in \textbf{biology}, reliable \emph{organismal development}; 
in \textbf{LLMs}, reasoning \emph{consistency} and \emph{alignment} across tasks. 
Crucially, this structure constrains \textbf{corpus-dependent variation} within \emph{interpretable latent geometry} -- preventing arbitrary or adversarial drift \citep{geva2021transformer, cobbe2021training}. 

\vspace{-2mm}
\begin{quote}
\emph{\textbf{What DNA is to the unity of multicellular life}, nDNA is to the coherence of LLM reasoning: a \textbf{stabilizing universal code} that enables \textbf{structured functional variation}.}
\end{quote}


\subsection*{Evolutionary and learning dynamics: convergence of principles}
Both \textbf{DNA} and \textbf{nDNA} are shaped by \emph{selection processes}. 
In \textbf{biology}, the genome has evolved under millennia of selective pressure, with \textbf{regulatory networks} fine-tuned to ensure \emph{robust development} and \emph{adaptability} \citep{alon2006introduction, davidson2006gene}. 
In \textbf{LLMs}, pretraining operates as an \emph{evolutionary analogue}: \textbf{stochastic gradient descent (SGD)} over massive corpora selects latent priors that minimize expected loss across tasks, with \emph{fine-tuning akin to epigenetic adjustment} \citep{bommasani2021opportunities, pfeiffer2021adapterfusion}. 

\[
\underbrace{
\mathcal{L}_{\mathrm{pretrain}}(\theta) = \mathbb{E}_{(x,y)} \left[ -\log p_\theta(y|x) \right]
}_{\textbf{SGD as selection pressure}}
\]

This \textbf{evolutionary parallel} explains why both systems exhibit \emph{clarity through complexity}: \textbf{layered hierarchies}, \textbf{probabilistic pathways}, and \textbf{interpretable modularity}. 
Where \textbf{biological evolution} yields \emph{modular gene regulatory networks} that ensure context-sensitive expression \citep{alon2006introduction}, \textbf{LLM training} yields \emph{modular latent structures} -- such as \textbf{attention heads} and \textbf{adapter modules} -- that scaffold \emph{task-specific reasoning} \citep{geva2021transformer, pfeiffer2021adapterfusion}.


\subsection*{Why corpus dependence matters}
Far from a flaw, \textbf{corpus dependence in \emph{n}DNA} is the signature of a \emph{flexible}, \emph{adaptive reasoning architecture}. 
Just as biological systems rely on \textbf{tissue-specific gene expression} to produce functional diversity from a \emph{universal genome} \citep{davidson2006gene, alon2006introduction}, large language models (LLMs) leverage \textbf{corpus-dependent latent scaffolding} to generate reasoning structures attuned to task demands, mirroring the reproducibility logic of biological variability quantification \citep{marioni2011rna}. 
By examining nDNA’s \textbf{spectral curvature} ($\kappa$), \textbf{thermodynamic length} ($\mathcal{L}$), and \textbf{belief vector norm} ($\|\mathbf{v}_\ell^{(c)}\|$), we gain a \textbf{diagnostic lens} for alignment, generalization, and safety \citep{belkin2019reconciling, still2012thermodynamic, olah2020zoom}:
\[
\mathcal{S}_{\mathrm{nDNA}}(c) = \int_{\gamma_c} \left( \alpha \kappa + \beta \mathcal{L} + \gamma \|\mathbf{v}_\ell^{(c)}\| \right) ds
\]
where $\gamma_c$ is the latent trajectory for corpus $c$. This latent geometry echoes Waddington’s epigenetic landscape where paths represent developmental fates \citep{waddington1957strategy}. 
Figure ~\ref{fig:ndna_groups} -- \textbf{QA tasks} evoke compact low-curvature paths (e.g.\ $\kappa \sim 0.012$--$0.03$, $\mathcal{L} \sim 0.47$--$0.53$) \citep{rajpurkar2016squad, kwiatkowski2019natural, joshi2017triviaqa}, while \textbf{reasoning tasks} elicit broader high-curvature paths (e.g.\ $\kappa \sim 0.005$--$0.04$) \citep{cobbe2021training, patel2021nlp, geva2021transformer}. 
\textbf{Dialogue corpora} produce shallow clustered scaffolds \citep{budzianowski2018multiwoz, li2016persona, zhang2018personalizing}; \textbf{commonsense tasks} yield oscillatory paths \citep{sap2019socialiqa, zellers2019hellaswag, talmor2019commonsenseqa}. 
nDNA aligns with interpretable AI goals \citep{zhang2018interpretable} and geometric decoding approaches \citep{narayanan2021decoding}.


This corpus dependence is \emph{not arbitrary noise} -- it reflects the model’s \textbf{learned latent regulatory logic}, analogous to the combinatorial control of \textbf{gene regulatory networks} that ensures \emph{context-sensitive yet robust gene expression} \citep{alon2006introduction, lewin2013genes}. 
Just as \emph{developmental disorders} arise when regulatory circuits misfire \citep{davidson2006gene}, misalignment or hallucination in LLMs can be traced to \emph{latent trajectories that diverge from expected scaffolding}. 
\textbf{nDNA analysis}, therefore, does not merely characterize model geometry -- it offers a \textbf{tool for interpretability, failure detection, and safe alignment}.

\vspace{-1mm}
\begin{quote}
\emph{\textbf{Corpus dependence in nDNA is the expression of reasoning plasticity}, bounded by universal latent priors much like gene networks balance flexibility with functional coherence.}
\end{quote}
\vspace{-1mm}

Moreover, the \textbf{universality of nDNA’s foundational structure} -- its pretrained manifold, architectural symmetries, and core alignment priors -- provides the \emph{stabilizing grammar} that constrains corpus-specific scaffolds within meaningful reasoning spaces \citep{vaswani2017attention, bommasani2021opportunities}. 
This is the latent equivalent of biology’s \textbf{genetic code} and \textbf{conserved transcriptional machinery}: an \emph{invariant substrate} that supports functional diversity without sacrificing coherence. 
By quantifying how nDNA paths \emph{bend}, \emph{stretch}, or \emph{steer} in response to task demands, we can map the model’s \textbf{cognitive landscape} -- and determine when it traces \emph{human-aligned reasoning} or drifts into failure modes.

\vspace{-1mm}
\begin{quote}
\emph{\textbf{What the genome is to life's functional unity}, nDNA is to the model’s reasoning coherence: a universal code that binds diversity into stability, and complexity into interpretability.}
\end{quote}
\vspace{-1mm}






\begin{table}[H]
\centering
\caption*{\centering \textbf{Mathematical comparison of DNA and nDNA structural layers}}
\begin{tabular}{p{3cm} | p{5.6cm} | p{5.6cm}}
\toprule
\textbf{Layer} & \textbf{DNA (Biology)} & \textbf{nDNA (LLM)} \\
\midrule
\textbf{Universal code} & Codon mapping $\phi: \Sigma^3 \to \mathcal{A}$, kernel $\neq \emptyset$, redundancy ensures error tolerance \citep{lewin2013genes} & Pretrained latent manifold; symmetries $G_{\mathrm{LLM}} \subset \mathrm{Aut}(V)$; generalization via equivariance \citep{bommasani2021opportunities} \\
\textbf{Context regulator} & Conditional $P(\text{gene ON}|\text{TF, epi})$; Bayesian gene networks \citep{alon2006introduction} & Conditional latent path $P(h_1,\dots,h_L|x)$; stochastic latent dynamics \citep{geva2021transformer} \\
\textbf{Path geometry} & Minimal energy path $\gamma^*$ in epigenetic landscape: $\int_\gamma \|\nabla V\| ds$ \citep{waddington1957strategy} & Latent geodesic minimizing cost: $\int_\gamma \|\nabla_\theta \log p(y|x)\|^2 ds$ \citep{still2012thermodynamic} \\
\textbf{Output mapping} & Fiber bundle: $\pi: E_{\mathrm{gene}} \to B_{\mathrm{cell}}$ & Fiber bundle: $\pi: E_{\mathrm{latent}} \to B_{\mathrm{task}}$ \\
\bottomrule
\end{tabular}
\end{table}





\begin{figure*}[t]
  \centering
  \begin{minipage}[t]{0.48\textwidth}
    \centering
    \includegraphics[width=\textwidth]{corpus_ndna/ndna_llama_qa_group.png}
    \subcaption{\textbf{QA group nDNA trajectories}: $\kappa$ ranges $\sim 0.012$--$0.03$, $L$ $\sim 0.47$--$0.53$, $\tau$ $\sim 0.006$--$0.014$. Trajectories are compact and consistently shaped across datasets, reflecting \textbf{shared task structure}.}
  \end{minipage}
  \hfill
  \begin{minipage}[t]{0.48\textwidth}
    \centering
    \includegraphics[width=\textwidth]{corpus_ndna/ndna_llama_dialogue_group.png}
    \subcaption{\textbf{Dialogue group nDNA trajectories}: $\kappa$ ranges $\sim 0.01$--$0.03$, $L$ $\sim 0.47$--$0.53$, $\tau$ $\sim 0.006$--$0.014$. Trajectories are shallow and tightly clustered, reflecting \textbf{low latent complexity} typical of conversational flow.}
  \end{minipage}
  \vspace{0.4cm}
  \begin{minipage}[t]{0.48\textwidth}
    \centering
    \includegraphics[width=\textwidth]{corpus_ndna/ndna_llama_reasoning_group.png}
    \subcaption{\textbf{Reasoning group nDNA trajectories}: $\kappa$ ranges $\sim 0.005$--$0.04$, $L$ $\sim 0.44$--$0.56$, $\tau$ $\sim 0.002$--$0.018$. Trajectories show \textbf{greater spread and complexity}, reflecting multi-step reasoning scaffolding.}
  \end{minipage}
  \hfill
  \begin{minipage}[t]{0.48\textwidth}
    \centering
    \includegraphics[width=\textwidth]{corpus_ndna/ndna_llama_commonsense_group.png}
    \subcaption{\textbf{Commonsense group nDNA trajectories}: $\kappa$ ranges $\sim 0.00$--$0.04$, $L$ $\sim 0.44$--$0.54$, $\tau$ $\sim 0.004$--$0.018$. Trajectories are intermediate in complexity, reflecting \textbf{varied latent demands of commonsense reasoning}.}
  \end{minipage}
  \vspace{-6mm}
  \caption{
    \textbf{nDNA trajectories across LLaMA vs. task groups.} 
    Each subplot visualizes \textbf{spectral curvature} ($\kappa_\ell$), \textbf{thermodynamic length} ($\mathcal{L}_\ell$), and \textbf{belief vector norm} ($\|\mathbf{v}_\ell^{(c)}\|$) layer-wise trajectories for representative datasets. 
    The structured variation illustrates that \emph{corpus dependence in nDNA is meaningful and interpretable}, reflecting task complexity rather than random noise. 
    \textbf{QA} and \textbf{dialogue} tasks activate \textbf{compact, smooth latent scaffolds with low curvature and modest belief steering}; \textbf{reasoning} tasks exhibit broader, more intricate geometry, with \textbf{increasing curvature, longer latent length, and stronger belief vector dynamics}. 
    \textbf{Commonsense} tasks show intermediate complexity with \textbf{oscillatory scaffolding, reflecting ambiguity and contextual switching}. 
    This figure demonstrates the core takeaway of our section: 
    \emph{like biological DNA, nDNA expresses differently in context, but remains bound by universal latent priors that ensure coherence, generalization, and alignment.}
    }
  \label{fig:ndna_groups}
\end{figure*}